The capacity of molnupiravir-mutagenized SARS-CoV-2 genomes to seed descendant clades constitutes a pivotal element of the thesis on the broken molecular evolutionary phylogenetic clock. Once the directional mutational process escapes the confines of a single treated host and enters the transmission chain, the pharmacological signature ceases to be a transient within-host artifact and becomes a heritable feature of circulating lineages. What follows is a continuous, mechanistically grounded account of how that seeding occurs, why it is phylogenetically detectable, and how it systematically undermines clock-like assumptions across the global tree.

Molnupiravir acts by elevating the error rate of the viral RNA-dependent RNA polymerase through incorporation of NHC-TP. The resulting genomes accumulate a heavy load of G-to-A and C-to-U transitions—often dozens to more than a hundred substitutions relative to the parental consensus—within a matter of days. Under the classical error-catastrophe model, the majority of such hypermutated genomes are expected to be non-viable. Empirical observation, however, demonstrates that a measurable fraction retain sufficient replicative fitness to reach detectable titers and to be shed. Clinical sequencing of treated patients reveals elevated transition-to-transversion ratios and the characteristic mutational spectrum while viral RNA remains quantifiable; in a subset of cases the same genomes persist beyond the treatment window. The survival of these variants is not anomalous once the stochastic nature of mutagenesis is considered: lethal mutations are frequent, yet genomes that by chance avoid critical catalytic or structural residues can continue to package and transmit. The elevated mutation supply simultaneously generates a broader cloud of minor variants, increasing the probability that at least one viable, signature-bearing genome is available for onward passage.

Phylogenetic reconstruction of global sequence databases supplies direct evidence that this passage has occurred repeatedly. Systematic surveys identify a discrete class of long terminal or near-terminal branches distinguished by an extreme enrichment of G-to-A and C-to-T substitutions. These branches appear almost exclusively after the clinical rollout of molnupiravir in late 2021 and early 2022, concentrate geographically in jurisdictions with high prescription volumes, and correlate demographically with age groups most frequently treated. Critically, a non-negligible proportion of these long branches are not phylogenetic dead-ends. Instead, they give rise to small but coherent descendant clusters whose subsequent mutations continue to accumulate under ordinary replication yet retain the ancestral molnupiravir spectral imprint. Documented examples include transmission chains recovered in the United Kingdom, the United States, Japan, New Zealand, Slovakia, Denmark, South Korea and Vietnam. In several instances the temporal interval between the inferred mutagenic event and the sampling of multiple secondary cases spans only weeks, demonstrating that the hypermutated genome can found a propagating lineage before purifying selection or further drift erases the diagnostic excess of transitions.

The seeding process therefore inserts abrupt, non-clock-like increments of genetic distance into the phylogeny. Under a conventional molecular-clock model, a long branch implies either an extended period of unobserved circulation or an elevated substitution rate sustained over many generations. When the branch is instead the product of a few days of intensive pharmacological mutagenesis followed by immediate transmission, both interpretations fail. The genetic distance is real, yet the calendar time required to generate it is radically compressed. Descendant clades inherit that compressed distance as a fixed offset; subsequent evolution proceeds from an already distorted baseline. Root-to-tip regressions and Bayesian dating analyses that treat all substitutions as products of a roughly constant endogenous rate consequently mis-estimate divergence times, inflate apparent evolutionary rates along affected lineages, and introduce heterogeneity that cannot be explained by host adaptation or immune selection alone.

Within the computational thesis this empirical pattern is formalized by coupling the directional mutagenesis module to multi-generational dynamics. Year-specific parameter matrices encode the progressive elevation of the mutagenic rate, the dominance of the G-to-A / C-to-U spectrum, and the concurrent decline in codon-usage bias. Each simulation step generates a mutational load whose composition matches the observed pharmacological signature; genomes that retain viability under the elevated load continue to update titer, fitness, glycosylation state and tropism. Because the model advances through successive generations (and, at annual intervals, through successive year matrices), the signature-bearing states are not reset. They propagate forward, interacting with the declining CUB, rising brain-tropism and exhaustion attractor terms. The numerical trajectory thereby illustrates how a single mutagenic pulse can seed a persistent dynamical regime whose statistical properties—accelerated observed mutation counts, spectral bias, and progressive phenotypic drift—diverge from those projected under a 2019-calibrated clock. In this sense the simulation does not merely reproduce isolated long branches; it demonstrates the conditions under which those branches can found self-sustaining descendant trajectories that continue to distort the phylogenetic signal for years.

The cumulative effect on the global phylogeny is a progressive contamination of the temporal scaffold. Pre-2022 calibration points remain relatively clean; post-2022 nodes become increasingly heterogeneous mixtures of ordinary replication and intermittent pharmacological jumps. Any attempt to extrapolate molecular-clock dates backward across the 2022 threshold—whether to refine the timing of early human emergence or to test the feasibility of multi-host cascades inside a narrow post-sampling window—must therefore confront the possibility that some of the observed distance is non-biological in origin. Conversely, forward projections that assume continued clock-like behavior systematically under-estimate the mutational supply available for adaptation once signature-bearing lineages are already circulating. The seeding of descendant clades is thus not a peripheral curiosity; it is the mechanism that converts a localized antiviral side-effect into a systemic perturbation of the evolutionary record.

In summary, the thesis holds that viable molnupiravir-mutagenized genomes routinely exit treated hosts, found secondary transmission chains, and thereby embed irreversible, directionally biased increments of genetic distance into the SARS-CoV-2 phylogeny. The resulting long branches and their descendant clades constitute empirical signatures of a broken molecular clock: substitutions that no longer map reliably onto calendar time, rate heterogeneities that track prescription geography rather than intrinsic biology, and a progressive erosion of the temporal calibration on which phylogenetic inference has historically relied.
